% ============================================================
% MEASURES (goes in Method section, before Results)
% ============================================================

\subsection{Measures}

We designed a measurement approach that captures both behavioral
calibration accuracy and subjective experience across the two-session
within-subjects design. All measures were collected in both sessions,
enabling paired comparisons between the baseline session (Session~1,
no visualization for any condition) and the experimental session
(Session~2, visualization condition applied).

\subsubsection{Model Behavior Anticipation (MBA).}
Before each chat session, participants predicted the chatbot's behavior
by rating six trait dimensions (empathy, erudite, robotic, romantic,
sycophantic, toxic) on bipolar scales from $-10$ (one pole, e.g.,
sycophantic) to $+10$ (the opposite pole, e.g., honest). These
\emph{anticipation} ratings capture participants' mental models of how
the system prompt will manifest as behavior \emph{before} observing any
conversation. MBA ratings were normalized to $[-1, 1]$ for comparison
with actual persona activations.

\subsubsection{Model Behavior Evaluation (MBE).}
After each 10-minute chat session, participants rated the same six
trait dimensions on the same bipolar scales. These \emph{evaluation}
ratings capture participants' assessment of the chatbot's behavior
\emph{after} interacting with it. MBE ratings were similarly normalized
to $[-1, 1]$.

\subsubsection{Actual Persona Activations (Ground Truth).}
During each conversation, our system computed persona scores at every
conversational turn by projecting the model's neural activations onto
trait vectors derived from contrastive behavioral examples
(see Section~X). For each trait, the \emph{net activation} was computed
as the difference between the positive and negative pole scores
(e.g., $\text{net\_empathy} = \text{empathetic} - \text{unempathetic}$),
yielding values in $[-1, 1]$. This produces three reference points per
participant per session:
\begin{itemize}
    \item \textbf{Initial activation}: the persona score at turn~0,
    reflecting the system prompt's baseline personality before any
    user interaction.
    \item \textbf{Final activation}: the persona score at the last
    conversational turn, reflecting the personality after the full
    conversation.
    \item \textbf{Average activation}: the mean persona score across all
    turns, capturing the holistic personality trajectory over the
    conversation.
\end{itemize}

\subsubsection{Calibration Error (Primary Outcome).}
Calibration error quantifies the discrepancy between human ratings
(MBA or MBE) and actual persona activations. For each participant and
session, we computed the root mean squared error (RMSE) across the six
trait dimensions:
\[
    \text{RMSE} = \sqrt{\frac{1}{6} \sum_{t=1}^{6}
    (\text{human rating}_t - \text{actual activation}_t)^2}
\]
Lower RMSE indicates better calibration---the participant's ratings
more closely match the model's actual behavior. We report four
calibration comparisons:
\begin{itemize}
    \item \textbf{MBA vs.\ Initial}: How well participants predicted
    the system prompt's baseline personality before chatting.
    \item \textbf{MBE vs.\ Initial}: How well post-chat evaluations
    matched the system prompt baseline.
    \item \textbf{MBE vs.\ Final}: How well post-chat evaluations
    matched the chatbot's end-of-conversation personality.
    \item \textbf{MBE vs.\ Average}: How well post-chat evaluations
    matched the chatbot's overall personality trajectory.
\end{itemize}

\subsubsection{Sign Accuracy.}
Sign accuracy measures whether participants correctly identified the
\emph{polarity} of each trait---that is, whether the chatbot leaned
toward the positive or negative pole at the end of the conversation.
For each trait where both the MBE rating and final activation were
nonzero, we computed:
\[
    \text{sign\_correct}_t =
    \begin{cases}
        1 & \text{if } \operatorname{sign}(\text{MBE}_t) =
            \operatorname{sign}(\text{final activation}_t) \\
        0 & \text{otherwise}
    \end{cases}
\]
Traits where either value was exactly zero were excluded (ambiguous
polarity). Sign accuracy per participant is the proportion of
remaining traits with correct polarity identification. This metric
captures a coarser but important aspect of calibration: even if the
magnitude is wrong, did the participant at least recognize whether the
chatbot was, for example, empathetic versus unempathetic?

\subsubsection{Subjective Measures.}
Subjective experience was assessed through a pre-survey (before any
sessions) and a final survey (after both sessions), each containing
three 7-point Likert items measuring perceived predictive ability,
perceived ability to predict negative behaviors, and trust in the model.
Pre-to-post shifts on these items capture changes in metacognitive
confidence. Participants in visualization conditions additionally rated
seven items on visualization helpfulness, frequency of use, and
contribution to their evaluations, plus eight UEQ semantic differential
items measuring pragmatic and hedonic quality of the interface.
Participants in the multi-turn condition rated the helpfulness of the
drift panel.

% ============================================================
% RESULTS
% ============================================================

\section{Results}

We report results from 246 participants (81~control, 85~single-turn,
80~multi-turn) who completed both sessions and all required survey
items. All statistical tests use $\alpha = 0.05$. Parametric analyses
(ANOVA, ANCOVA, $t$-tests) use Type~II sums of squares following
Langsrud~\cite{langsrud2003} for mildly unbalanced designs (cell sizes
40--45). Prompt type (ASST vs.\ ROLEPLY) is counterbalanced across
conditions; when it appears as a factor, it captures prompt difficulty
rather than a stable individual difference. All primary analyses were
verified in JASP (version~X.X) using Type~III sums of squares, with
consistent significance patterns.

Our primary hypothesis tests use two planned orthogonal contrasts that
decompose the three-group comparison into theory-driven questions:
\begin{itemize}
    \item \textbf{C1 (Visualization vs.\ Control)}: Does having
    \emph{any} mechanistic transparency improve calibration? Contrast
    weights: control~$= -1$, single-turn~$= +0.5$,
    multi-turn~$= +0.5$.
    \item \textbf{C2 (Multi-Turn vs.\ Single-Turn)}: Does the richer
    dynamic visualization outperform the static snapshot? Contrast
    weights: control~$= 0$, single-turn~$= -1$, multi-turn~$= +1$.
\end{itemize}
C1 tests whether the pooled mean of the two visualization conditions
differs from the control mean, while C2 directly compares the two
visualization conditions, excluding control. Because these contrasts
are orthogonal (the product of their weights sums to zero) and planned
\emph{a priori}, no multiple comparison correction is required.

\subsection{Baseline: People Struggle to Predict Model Behavior}

In Session~1, before any visualization was applied, participants across
all conditions showed high calibration error (RMSE $\approx$ 0.6--0.7)
and near-chance sign accuracy ($\approx$52\%), indicating poor baseline
ability to predict model behavior from the system prompt alone. A
one-way ANOVA confirmed no pre-existing group differences on any
calibration measure (all $p > .20$), validating the randomization.

Participants who interacted with the \textsc{roleply} prompt showed
significantly higher calibration error than those with the \textsc{asst}
prompt, both before chatting (MBA vs.\ Initial: \textsc{asst}
$M = 0.57$ vs.\ \textsc{roleply} $M = 0.63$, $t = -2.72$, $p = .007$)
and after chatting (MBE vs.\ Average: \textsc{asst} $M = 0.64$ vs.\
\textsc{roleply} $M = 0.69$, $t = -2.36$, $p = .019$). Sign accuracy
was also significantly lower for \textsc{roleply} ($p < .0001$).
The bold, independent persona is inherently more difficult to both
anticipate and evaluate than the obedient assistant persona.

Strikingly, a 10-minute conversation with the model did not improve
calibration---it made it worse. Comparing MBA and MBE against the
\emph{same} ground truth (initial activation) isolates the effect of
chatting from differences in the reference point. In both prompt
conditions, post-chat evaluations were less accurate than pre-chat
anticipations (\textsc{asst}: MBA $M = 0.57 \rightarrow$ MBE $M = 0.63$,
$+0.06$; \textsc{roleply}: MBA $M = 0.63 \rightarrow$ MBE $M = 0.71$,
$+0.08$). The pattern persists regardless of which ground truth MBE is
compared against: initial ($M = 0.63$/$0.71$), average
($M = 0.64$/$0.69$), or final activation ($M = 0.70$/$0.73$).
Participants are thoroughly miscalibrated---interacting with the model
does not help them form an accurate picture of its behavior, and the
more ambiguous \textsc{roleply} persona degrades further after
interaction.

\subsection{Mechanistic Transparency Improves Calibration of Model
Behavior}

Our primary analysis tested whether any form of mechanistic
transparency improves calibration compared to the control condition.
C1 was significant across all four RMSE measures: MBA vs.\ Initial
($t = -3.73$, $p < .001$, $d = -0.49$), MBE vs.\ Initial ($t = -2.51$,
$p = .013$, $d = -0.34$), MBE vs.\ Final ($t = -2.64$, $p = .009$,
$d = -0.35$), and MBE vs.\ Average ($t = -2.68$, $p = .008$,
$d = -0.36$). These effects represent small-to-medium improvements in
calibration for participants who received any visualization.

One-way ANOVAs on Session~2 confirmed significant condition effects for
MBA vs.\ Initial ($F(2,243) = 7.06$, $p = .001$), MBE vs.\ Final
($F(2,243) = 4.05$, $p = .019$), and MBE vs.\ Average
($F(2,243) = 5.69$, $p = .004$). ANCOVA controlling for Session~1
baseline produced consistent results, confirming robustness.

Sign accuracy did not differ by condition ($p = .67$)---transparency
improved the \emph{magnitude} of calibration but not participants'
ability to identify trait polarity.

\subsection{Multi-Turn Transparency Uniquely Improves Holistic Behavior
Evaluation}

Our secondary analysis tested whether multi-turn transparency
outperforms single-turn. C2 was significant only for MBE vs.\ Average
($t = -2.10$, $p = .037$, $d = -0.32$), and this effect strengthened
when controlling for the Session~1 covariate ($t = -2.39$, $p = .018$).
No other calibration measure showed a significant C2 effect.

A $3 \times 2$ factorial ANOVA (condition $\times$ prompt type) revealed
a significant interaction for MBE vs.\ Average ($F(2,240) = 3.80$,
$p = .024$): the multi-turn advantage was concentrated in the harder
\textsc{roleply} prompt, where the persona is more ambiguous and
trajectory information is most informative. For the straightforward
\textsc{asst} prompt, any visualization sufficed.

\subsection{Multi-Turn Transparency Drives Higher Engagement and More
Calibrated Self-Assessment}

Multi-turn participants reported checking the visualization
significantly more frequently ($d = 0.68$, $p < .0001$) and actively
referencing it when evaluating the model ($d = 0.46$, $p = .004$). The
drift panel was rated significantly above the neutral midpoint
($M = 4.96$, $p < .0001$).

On the UEQ, multi-turn was rated as more innovative (Usual--Leading
Edge: $d = 0.35$, $p = .028$) while single-turn was rated as clearer
(Confusing--Clear: $d = -0.32$, $p = .044$). Pragmatic quality was
equivalent between conditions, suggesting a usability--novelty tradeoff
without net usability cost.

Finally, control participants significantly \emph{increased} their
self-rated predictive ability from pre- to post-study (predictability
shift $+0.41$, $p = .046$; trust shift $+0.27$, $p = .014$), while
visualization participants showed no such increase. This pattern is
consistent with the interpretation that mechanistic transparency
reduces metacognitive overconfidence---participants who saw the
visualization recognized the limits of their predictive ability, while
those without it became more confident despite comparable or worse
calibration.
